% RESEARCH_DRAFT_SHA256: 6951da8dee10153239dda9fa472a2d71e269efd595b241da3f4858e18daeb681
% Layout and source-reference wrappers expanded for the web renderer.
\documentclass[11pt]{article}
\usepackage[margin=25mm]{geometry}
\usepackage{fontspec,unicode-math}
\setmainfont{Latin Modern Roman}
\setmathfont{Latin Modern Math}
\usepackage{amsmath,amssymb,mathtools,xurl,hyperref}
\hypersetup{colorlinks=true,urlcolor=blue,linkcolor=blue}
\setcounter{secnumdepth}{0}
\setlength{\parindent}{0pt}
\setlength{\parskip}{5pt}
\setlength{\emergencystretch}{3em}
\newcommand{\sref}[2]{\hyperlink{source:#1:#2}{[S#2]}}
\newcommand{\cataloguescope}[1]{\paragraph{Source record.}#1}
\newcommand{\floor}[1]{\left\lfloor#1\right\rfloor}
\begin{document}
\section*{Research-attempt provenance}
This standalone research attempt retains the complete original definition body below,
including its historical source and preparation notes. Those notes describe the
earlier catalogue preparation; the new attempt and its verification are identified
separately after that body. No original ID, catalogue status, or source statement is changed.

\section{Reducing the Exponential Tower in Constructive Degree Bounds}\label{statement}

\subsection{Definitions and mathematical statement}

\paragraph{Shared notation.}Standard number systems are $\mathbb N,\mathbb Z,\mathbb Q,\mathbb R,\mathbb C$. Specialized notation and index ranges follow the cited source and the definitions below.


A tower of exponentials iterates exponentiation: for example $2^{2^{2^x}}$ has three displayed exponential stages. A constructive degree bound limits the polynomial degrees occurring in the specified construction.
\paragraph{Record scope and qualifications.}
The catalogue title is retained as metadata. Named certificates, contracts, bounds or frameworks in that title are not asserted as results of the original source extract.
The construction, degree parameters, and existing bound are absent from the source extraction; no missing bound is invented.

Source-record limitation: The degree bound, parameters, and mathematical object to which the unnamed five-exponential tower applies are absent.
\paragraph{Statement recorded in the supplied source.}
Is the tower of five exponentials in the constructive degree bound necessary, or can its height be substantially reduced? \href{https://doi.org/10.4171/owr/2017/14}{[S2]}
\paragraph{Source status.}
The supplied catalogue labels this source record open; that metadata was checked against the supplied snapshot on 2026-10-11. The label is preserved for provenance and is not an independent verification of present open status or mathematical correctness. The catalogue literature assessment is dated 2026-08-21.
\subsection{Short English statement}
Reduce the height of the cited degree-bound exponential tower.
\subsection{Sources}
\begin{itemize}
\item[S1] Ulam AI and the UnsolvedMath Contributors, UnsolvedMath: A Curated Collection of Open Mathematics Problems, catalogue entry. Catalogue curation and metadata: CC BY 4.0. The source extract is retained; definitions, English summaries and qualifications are editorial additions.
\url{https://huggingface.co/datasets/ulamai/UnsolvedMath}.
\item[S2] Real Algebraic Geometry With a View Toward Moment Problems and Optimization (2018), Oberwolfach Reports, DOI 10.4171/owr/2017/14, pp. 771--862
\url{https://doi.org/10.4171/owr/2017/14}.
\item[S3] Lombardi--Perrucci--Roy, elementary-recursive bound
\url{https://arxiv.org/abs/1404.2338}.
\item Source-specific Oberwolfach Reports licences, including CC BY-SA 4.0 where applicable, are retained; see THIRD\_PARTY\_NOTICES.md.
\end{itemize}
\paragraph{Preparation and review.}
The original source and UnsolvedMath contributors retain attribution for the source extract. Codex (GPT-6; exact serving revision not exposed) assisted with definitions, English summaries, scope notes and formatting. No new human mathematical authorship or independently verified solution is asserted. Catalogue identity checked against the complete supplied registry and current repository. Distinct upstream records, including repeated questions and broad or damaged extracts, are retained for complete coverage of the source CS category. A repeated record is not a new independent conjecture. Complete mathematical-equivalence and current-literature review has not been performed. Preparation and final review follow contribution prompts C and D, with source gaps disclosed. Structural and rendering checks do not establish a new theorem or open status.
\label{end:statement}

\section{Literature and context check}
The primary OWR 14/2017 talk, pp.~779--781, was read. The displayed question concerns effective Positivstellensatz and Hilbert's seventeenth problem, not an unspecified tower in a different algorithm. The abstract of Lombardi--Perrucci--Roy, arXiv:1404.2338v3, was also consulted and explicitly reports a tower of five exponentials. The very long constructive proof and later improvements have not been checked here. The calculations below explain why two tempting reductions do not settle the general degree-bound question.
\section{Research attempt: Low-dimensional certificates and the denominator obstruction}

\paragraph{Univariate factorization gives a small certificate over the reals.}
Let $p\in\mathbb R[t]$ be nonnegative on $\mathbb R$. If $p=0$, its zero-square representation suffices. Otherwise its degree is even, say $2d$, its leading coefficient is positive, and each real root has even multiplicity: an odd multiplicity would change the sign across that root. Using the fundamental theorem of algebra, factor the remaining nonreal roots into conjugate pairs. Thus $p$ is a positive constant times even powers of real linear factors and positive quadratics of the form
\[
 (t-a)^2+b^2,\qquad b>0.
\]
Each factor is a sum of two squares, and
\[
 (u^2+v^2)(r^2+s^2)=(ur-vs)^2+(us+vr)^2
\]
combines them. Induction gives $p=u^2+v^2$ with $u,v\in\mathbb R[t]$ of degrees at most $d$. The degree bound follows because each multiplication adds at most half the degree of the multiplied factor. It is sharp in this representation class for $p=t^{2d}$: squares of polynomials of degree less than $d$ cannot have degree $2d$. This is a known elementary univariate argument. It uses real factor coefficients; for rational input it does not by itself produce a rational-coefficient certificate over the original ordered coefficient field.

\paragraph{The polynomial-SOS shortcut fails in two variables.}
Consider the Motzkin polynomial
\[
 M(x,y)=x^4y^2+x^2y^4+1-3x^2y^2.
\]
It is nonnegative: for nonnegative $a,b,c$, $a+b+c\ge3(abc)^{1/3}$; applying this to $x^4y^2,x^2y^4,1$ gives the claim. One proof of this elementary inequality is to minimize $a+b+c$ at fixed positive product: replacing two unequal entries by their geometric mean decreases their sum, and a minimizer exists after restricting to a bounded sublevel set, so the minimum occurs when all three agree. A zero product is immediate.

Suppose $M=\sum_j p_j^2$ with real polynomials. Every $p_j$ has degree at most three: otherwise the highest homogeneous terms give a nonzero sum of squares of higher degree, which cannot cancel. The coefficient of $x^6$ forces every $x^3$ coefficient in $p_j$ to vanish, and similarly the coefficient of $y^6$ removes every $y^3$ term. Then the zero coefficient of $x^4$ forces every $x^2$ coefficient to vanish, and similarly every $y^2$ coefficient. The zero coefficients of $x^2$ and $y^2$ now force every $x$ and $y$ coefficient to vanish. The only possible remaining monomials in $p_j$ are $1,xy,x^2y,xy^2$. Their squares have $x^2y^2$ coefficient equal to the sum of the squares of the $xy$ coefficients. This is nonnegative, contrary to the coefficient $-3$ in $M$. Thus a general method cannot simply replace rational-function squares by polynomial squares.

\paragraph{An explicit denominator repair, with its actual degree.}
Set $s=x^2+y^2$ and
\[
 F_1=xy(s-2),\qquad F_2=x(y^2-1),\qquad F_3=y(x^2-1).
\]
Expansion gives the exact identity
\[
 sM=F_1^2+F_2^2+F_3^2.
\]
Consequently in the rational-function field $\mathbb R(x,y)$,
\[
 M=\sum_{j=1}^3\left[\left(\frac{xF_j}{s}\right)^2+
                              \left(\frac{yF_j}{s}\right)^2\right].
\]
The common denominator has degree two and the displayed numerators have degree at most five. This is an identity of rational functions, so the common zero of the denominator at the origin does not invalidate it; it is not a claim that each displayed quotient is defined everywhere as a real function. The example shows concretely how a failed polynomial certificate can be repaired, but supplies no control of denominator growth for arbitrary dimension and degree. Extending the univariate factorization proof to multivariate irreducible factors is precisely where the argument stops.
\section{Verification and adversarial audit}
The Motzkin coefficient argument checks all possible degree-at-most-three monomials, rather than inferring non-SOS from nonnegativity. The displayed denominator identity was expanded exactly using integer monomial dictionaries in Python; it matched coefficient by coefficient. Zero and constant univariate polynomials, coefficient-field restrictions and denominator zeros were checked. The fundamental theorem of algebra is the sole standard existence theorem used for the univariate factorization; no effective multivariate Positivstellensatz theorem is invoked in the attempted proof.
\paragraph{Audit corrections.} Kept real-coefficient univariate factorization separate from the original ordered-field construction. Rejected the invalid polynomial-SOS replacement and supplied an exact rational repair instead. No general tower improvement is inferred from low-dimensional examples.
\paragraph{External computation.} Python 3.14.2, exact integer coefficient dictionaries: checked $(x^2+y^2)M-F_1^2-F_2^2-F_3^2=0$; the resulting dictionary was empty. The companion script \texttt{check\_polynomial\_identity.py} records the full input and check.
\section{Scope and remaining gaps}
Neither the general lower bound nor a smaller-height constructive upper bound is proved. The source's ordered-field setting and its sign-condition certificates are broader than real univariate factorization or one two-variable example. The first missing step is a uniform multivariate construction controlling the degrees of denominators and identities; an individual small certificate does not control worst-case tower height.
\paragraph{Integration review.} Codex (GPT-6; exact serving revision not exposed) reviewed the written arguments and expanded layout and source-reference wrappers. Repository metadata and the production web renderer were checked using Python 3.13 and JavaScript with MathJax 3.2.2. This model-assisted review is not independent human proof verification; no PDF compilation or formal proof checking is asserted.

\section{Completion Estimate}
\noindent\textbf{COMPLETION: 0\%}
The calculations below the preserved source statement investigate a concrete route,
but establish no substantiated new progress toward the entire original target.
Known special cases and failed sufficient conditions are not credited as a solution.
\section{References}
The original source references and catalogue attribution are preserved in the
Sources section above. The following URLs were actually consulted for this attempt;
retention of an additional original reference does not claim fresh consultation of it.
\begin{itemize}
\item \url{https://ems.press/content/serial-article-files/46677} (consulted 2026-10-11). 
\item \url{https://arxiv.org/abs/1404.2338} (consulted 2026-10-11). 
\end{itemize}
\section{Model and contributor disclosure}
OpenAI gpt-6-astra, requested reasoning configuration ultra, generated this research
attempt on 2026-10-11. The exact serving revision is not exposed. Codex GPT-6
performed orchestration, repository-format review, and written-proof review of this attempt; its exact serving revision
is also not exposed. The mathematical route, boundary checks and final adversarial
self-audit followed contribution prompts A then B in
\texttt{docs/contribution\_prompts.md}, after reading \texttt{CONTRIBUTING.md}.
Python 3.14.2 assembled the standalone documents; any mathematical computation is
specified in Verification. No theorem prover was used. No human mathematical
authorship, substantial human proof correction, or independent human proof checking
is asserted. Model-assisted review, source retrieval, and structural checks do not
establish correctness or novelty. Full private reasoning transcripts are unavailable
as reproducibility artifacts; the complete arguments, inputs and outcomes actually
used are stated here or in the companion batch audit.
\section{Final status}
Unresolved as a new research contribution. This is a concrete route investigation,
not a claim of new complete resolution or a current-status certification. A known
answer reproduced for context does not become a new contribution. The catalogue's
source-status label is preserved; this attempt status does not assert that a
source-labelled solved question is open.
\end{document}
